\section{Related Work}\label{related_work}

Die Erkennung von Cyberangriffen in Internet of Things (IoT)-Netzwerken mittels Machine Learning (ML) ist ein aktives und wachsendes Forschungsfeld. Die folgenden Abschnitte gliedern die zehn verwendeten Quellen in vier thematische Cluster: Datensätze und Benchmarks, ML-basierte Ansätze zur Angriffserkennung, Behandlung von Klassenungleichgewicht sowie Erklärbarkeit von IDS-Modellen.

\subsection{Datensätze und Benchmarks für IoT-Sicherheit}
Die empirische Forschung im Bereich IoT-Sicherheit ist grundlegend auf realistische und repräsentative Benchmarkdatensätze angewiesen. Neto et al. \cite{neto2023ciciot2023} haben mit dem CICIoT2023-Datensatz einen umfassenden Benchmark entwickelt, der Netzwerkverkehr von 105 realen IoT-Geräten unter 33 verschiedenen Angriffsbedingungen enthält. Die Angriffe sind in sieben Kategorien organisiert: Distributed Denial of Service (DDoS), Denial of Service (DoS), Mirai, Reconnaissance, Spoofing, Brute-Force und Web-based Attacks. Ein zentrales Merkmal des Datensatzes ist, dass die Angriffe ausschließlich von IoT-Geräten gegen andere IoT-Geräte ausgeführt wurden, was eine besonders realitätsnahe Abbildung realer Bedrohungsszenarien ermöglicht \cite{neto2023ciciot2023}. Der CICIoT2023-Datensatz weist dabei ein ausgeprägtes Klassenungleichgewicht auf, das bei der Modellbewertung methodische Herausforderungen mit sich bringt und besondere Aufmerksamkeit erfordert.

\subsection{ML-basierte Ansätze zur Angriffserkennung in IoT-Netzwerken}
Hamedani et al. \cite{hamedani2025iotsurvey} liefern mit ihrer umfassenden Übersichtsarbeit einen systematischen Überblick über ML-basierte Sicherheitslösungen für IoT-Netzwerke aus den Jahren 2020 bis 2024. Die Autoren klassifizieren ML-Techniken nach ihrer spezifischen Anwendung in der IoT-Sicherheit und identifizieren offene Forschungsfragen, darunter insbesondere den Mangel an interpretierbaren Modellen für ressourcenbeschränkte Umgebungen \cite{hamedani2025iotsurvey}. Kikissagbe und Adda \cite{kikissagbe2024review} ergänzen diesen Überblick durch eine systematische Analyse von supervised, unsupervised und hybriden ML-Ansätzen für Intrusion Detection Systems (IDS) in IoT-Umgebungen. Ihre Arbeit zeigt, dass baumbasierte Ensemble-Methoden konsistent zu den leistungsstärksten Ansätzen gehören, gleichzeitig aber Fragen der Interpretierbarkeit und Ressourceneffizienz weitgehend unbeantwortet bleiben \cite{kikissagbe2024review}.

Auf dem CICIoT2023-Datensatz selbst haben Raturi et al. \cite{raturi2026exploratory} eine vergleichende Evaluation mehrerer ML-Algorithmen durchgeführt. Sie zeigen, dass baumbasierte Modelle und Ensemble-Methoden die höchsten Balanced-Accuracy-Werte erreichen, wobei Gradient Boosting in der Mehrklassenklassifikation den besten Wert unter den getesteten Modellen erzielt \cite{raturi2026exploratory}. Die Autoren führen einen zweistufigen hierarchischen Klassifikationsansatz ein, der das starke Klassenungleichgewicht des Datensatzes strukturell adressiert: In Stufe 1 wird zwischen DoS/DDoS-Verkehr und Nicht-DoS-Verkehr unterschieden, in Stufe 2 erfolgt die Mehrklassenklassifikation der verbleibenden Kategorien \cite{raturi2026exploratory}. Almahaqeri et al. \cite{almahaqeri2026gradient} evaluieren ebenfalls Gradient Boosting auf dem CICIoT2023-Datensatz und wählen eine andere Strategie zur Behandlung des Klassenungleichgewichts, nämlich stratifiziertes Undersampling. Alharby \cite{alharby2025ciciot} vergleicht auf demselben Datensatz mehrere ML-Modelle einschließlich Gradient Boosting und Decision Trees unter Einsatz von Principal Component Analysis (PCA) als Dimensionsreduktionstechnik und dokumentiert zusätzlich Trainings- und Vorhersagezeiten.

\subsection{Behandlung von Klassenungleichgewicht in IDS-Datensätzen}
Das Klassenungleichgewicht stellt eine der zentralen methodischen Herausforderungen bei der Entwicklung ML-basierter IDS dar. Hosseini et al. \cite{hosseini2025imbalance} untersuchen systematisch die Auswirkungen verschiedener Grade von Klassenungleichgewicht auf die Leistung von ML-Algorithmen in IDS-Datensätzen und vergleichen verschiedene Resampling-Strategien. Ihre Ergebnisse zeigen, dass die einfache Accuracy als Bewertungsmetrik bei unausgewogenen Datensätzen zu systematisch verzerrten Urteilen führt, da Modelle, die seltene Angriffsklassen vollständig ignorieren, rechnerisch hohe Accuracy-Werte erzielen können \cite{hosseini2025imbalance}. Raturi et al. \cite{raturi2026exploratory} greifen diese Problematik auf und schlagen die Balanced Accuracy als fairere Bewertungsgrundlage vor, definiert als arithmetisches Mittel aus Sensitivität und Spezifität. Sie belegen empirisch am Beispiel des Hidden Markov Models, dass ein Modell trotz eines Recall-Werts von 0,999 für DoS-Angriffe in der Mehrklassenklassifikation auf einen Recall-Wert von 0,13 einbrechen kann, was bedeutet, dass 87 Prozent aller Nicht-DoS-Angriffe unerkannt bleiben \cite{raturi2026exploratory}. Almahaqeri et al. \cite{almahaqeri2026gradient} verfolgen einen alternativen Ansatz und setzen stratifiziertes Undersampling ein, um das Klassenungleichgewicht zu reduzieren, anstatt eine hierarchische Klassifikationsarchitektur zu verwenden.

\subsection{Erklärbarkeit von ML-basierten IDS}
Der Einsatz von Methoden der erklärbaren Künstlichen Intelligenz (XAI) in IDS ist ein jüngeres, aber rasch wachsendes Forschungsfeld. Ogunseyi et al. \cite{ogunseyi2026xai_review} analysieren in einem PRISMA-basierten systematischen Review 129 peer-reviewte Studien aus den Jahren 2018 bis 2025 und stellen fest, dass SHapley Additive Explanations (SHAP) und Local Interpretable Model-agnostic Explanations (LIME) in 98 Prozent aller untersuchten XAI-IDS-Studien eingesetzt werden \cite{ogunseyi2026xai_review}. Ein zentrales Ergebnis ihrer Analyse ist, dass die Integration von post-hoc-XAI-Methoden die Erkennungsgenauigkeit der zugrundeliegenden Modelle nicht beeinträchtigt, da diese Methoden nach der Vorhersage angewendet werden, ohne die Modellparameter zu verändern \cite{ogunseyi2026xai_review}. Mohale und Obagbuwa \cite{mohale2025xai} evaluieren mehrere ML-Modelle, darunter Decision Trees und XGBoost, gemeinsam mit SHAP, LIME und ELI5 auf einem IDS-Datensatz und zeigen, dass SHAP die konsistentesten und zuverlässigsten Erklärungen liefert. Sie demonstrieren darüber hinaus, dass XAI-Methoden verborgene Schwächen im Trainingsprozess sichtbar machen können, die durch reine Accuracy-Bewertung verborgen bleiben \cite{mohale2025xai}. Hermosilla et al. \cite{hermosilla2025xai_forensic} vergleichen SHAP und LIME auf XGBoost- und TabNet-Modellen im Kontext der forensischen Intrusion Detection und führen drei quantitative Metriken zur Bewertung der Erklärungsqualität ein: Fidelität, Konsistenz und Jaccard-Ähnlichkeit. Ihre Ergebnisse zeigen, dass SHAP eine höhere Fidelität und Konsistenz erzielt als LIME, während beide Verfahren vergleichbare Jaccard-Ähnlichkeitswerte aufweisen \cite{hermosilla2025xai_forensic}.

\subsection{Identifizierte Forschungslücken und Positionierung der eigenen Arbeit}
Die Analyse der bestehenden Literatur offenbart eine konsistente und direkt belegbare Forschungslücken: Keine der identifizierten Studien vereint gleichzeitig alle drei folgenden Elemente auf dem CICIoT2023-Datensatz. Erstens die Verwendung der Balanced Accuracy als primäre Bewertungsmetrik. Zweitens die Integration von SHAP als post-hoc-Erklärungskomponente. Drittens eine systematische Ablation Study, die den Einfluss der PCA-Dimensionsreduktion auf Modellleistung und SHAP-Feature-Importance quantifiziert.

Raturi et al. \cite{raturi2026exploratory} verwenden zwar die Balanced Accuracy und den CICIoT2023-Datensatz, verzichten jedoch vollständig auf jede Form der Erklärbarkeitsanalyse. Almahaqeri et al. \cite{almahaqeri2026gradient} und Alharby \cite{alharby2025ciciot} arbeiten ebenfalls auf demselben Datensatz, integrieren aber weder Balanced Accuracy als primäre Metrik noch eine SHAP-Komponente. Mohale und Obagbuwa \cite{mohale2025xai} sowie Hermosilla et al. \cite{hermosilla2025xai_forensic} liefern fundierte SHAP-Integrationen mit quantitativer Qualitätsbewertung, arbeiten jedoch nicht auf dem CICIoT2023-Datensatz.

Die vorliegende Arbeit schließt diese Lücke, indem sie einen Gradient-Boosting-Klassifikator auf dem CICIoT2023-Datensatz mit Balanced Accuracy als primärer Bewertungsmetrik und einer vollständigen SHAP-Analyse auf beiden Klassifikationsstufen verbindet. Damit leistet sie einen eigenständigen Beitrag, der über die Summe der bestehenden Einzelansätze hinausgeht.
